\section{Related Work}
\label{sec:related_work}

\subsection{Building Graphs and Checking Extracted Facts}

Graph construction turns source data into linked facts \citep{ref_ji2021}.
DBpedia, YAGO, and Freebase build knowledge bases from structured and partly
structured sources \citep{ref_auer2007,ref_suchanek2007,ref_bollacker2008}.
Quality checks help select facts during this process.
Knowledge fusion combines evidence from several sources to score candidate
facts \citep{ref_dong2014}.
This connects extraction with the question of which facts to keep.

Learning methods use graph patterns to predict links
\citep{ref_bordes2013,ref_wang2017,ref_schlichtkrull2018}.
Surveys describe how LLMs use text to build, complete, and combine KGs
\citep{ref_pan2023,ref_pan2025,ref_bian2025}.
Reasoning prompts and prompt patterns guide model responses
\citep{ref_wei2023,ref_white2023}.
Retrieval-augmented generation adds source text to the model input
\citep{ref_lewis2020}.
Neo4j Labs' LLM Knowledge Graph Builder applies such models to documents
with user-set schemas \citep{ref_neo4j_llm_graph_builder}.
These studies provide ways to produce candidate facts.
Our task starts from an extracted graph and uses the source document to
check proposed field values.

\subsection{Graph Quality and Check Scopes}

Quality assessment asks whether stored facts are correct, complete,
consistent, and current \citep{ref_zaveri2016,ref_paulheim2017,ref_wang2021}.
One route checks how well the data fit observed patterns.
Statistical methods find unusual values from their distributions
\citep{ref_wienand2014,ref_paulheim2014}.
Embedding scores measure how well triples fit learned graph patterns
\citep{ref_nickel2016}.
Another route states the requirements for valid records through tests and
graph constraints \citep{ref_kontokostas2014,ref_shacl2017}.
Human review helps assess cases with unclear answers \citep{ref_acosta2013}.

These approaches turn quality goals into checks that locate different
kinds of errors.
For document graphs, local structure, field rules, and source evidence
provide linked views of the same record.
Our assessment profile brings repeated facts, field violations, and source
mismatches into one report.
This report connects the detected problems with repair choices.

\subsection{Repairing Graphs and Choosing Actions}

Repair involves recovering missing facts and preserving correct information.
Completion methods predict missing links from paths and embeddings
\citep{ref_lao2011,ref_bordes2013,ref_yang2015,ref_shi2017}.
Neural models also learn scores for possible links
\citep{ref_dettmers2018,ref_zhu2021}.
Path-based reinforcement learning supports query answering
\citep{ref_das2018}, while few-shot learning supports link prediction from
few examples \citep{ref_chen2019}.
These methods provide candidate facts for graph completion.

Rules give proposed facts a structure that can be checked.
Rule mining finds positive and negative patterns in knowledge bases
\citep{ref_galarraga2013,ref_ortona2018}.
Embedding training can include logical rules and simple constraints
\citep{ref_guo2016,ref_ding2018}.
RNNLogic and pLogicNet combine logical rules with learned predictions
\citep{ref_qu2021,ref_qu2019}.
These studies connect fact prediction with rules that guide learning and
checking.

Language models provide another route to repair through prompts and context.
Lin et al.\ use violated SHACL constraints and graph context in repair
prompts \citep{ref_lin2025}.
ReAct links model reasoning with actions that gather information
\citep{ref_yao2023react}.
We draw on these designs for the SHACL-context and ReAct-style baselines.
Building on these studies, we connect source evidence with field constraints
and edit choice.
The main path filters proposed triples, while a scale-aware selector
compares trial edits by quality gain and cost under graph bounds.
